\scanpage{63}
\mspage{12}
La donnée d'un $R_i^\omega$, \inserted{qu'on va écrire $R_{\omega i}^{\omega}$, équivaut à} celle d'un repère \inserted{$r=(i,\omega)$} de $J$, et définit \uline{deux} \deleted{\unsure{rég}} régions $\widetilde D_r$ qui le contiennent, correspondant aux deux repères de $(J,\unsure{\uline{\varepsilon}|\varepsilon})$ qui prolongent $r$ -- et il y a donc deux régions correspondantes $\widetilde\Sigma_r$, qui sont les seules qui contiennent $R_i^\omega$ comme \uline{point intérieur} (NB. dans la région $\widetilde\Sigma_{\unsure{r_\varepsilon|r_0}}$ standard, les deux autres points $R$ à part $R_0^+$, qui est intérieur, sont $R_0^-$ et $R_1^-$, qui ne le sont pas). Mais si la sphère $\Sigma$ est orientée i.e. si on s'est donné un élément marqué \inserted{$\unsure{\bar\omega|\varpi}_+$} de $\uline{\unsure{\bar\omega|\varpi}}=\uline{\omega}\wedge\uline{\varepsilon}$, ou encore un isom $\unsure{\uline{\varepsilon}|\uline{\omega}}\simeq\uline{\varepsilon}$, alors le point $R_i^\omega$ détermine un $\varepsilon$ \inserted{$\in\uline{\varepsilon}$} privilégié, par la condition que $\varepsilon$ corresponde à $\omega$ par $\uline{\omega}\simeq\uline{\varepsilon}$ ou encore que $\omega\wedge\varepsilon=\unsure{\bar\omega|\varpi}_+$, ce qui signifie aussi que l'orientation $P_+S_0^+R_0^+Q_\infty P_+$ de la région $\deleted{\widetilde\Sigma_r}\widetilde D_r$ est \deleted{celle de} $\unsure{\bar\omega|\varpi}_+$. Donc on peut alors identifier canoniquement, via $\unsure{\Psi_r|\psi_r}$, $\Pi_1(\widetilde\Sigma,R_i^\omega)$ à $\Pi_{0,3}$, par un iso (qui dépend du choix de l'orientation $\unsure{\bar\omega|\varpi}_+$ de $\Sigma$)
\[
(43),\qquad\psi_i^\omega:\Pi_{03}\xrightarrow{\sim}\Pi_1(\widetilde\Sigma,R_i^\omega)
\]
\margin{62}
